%% ma: L12-B13-0032
%% lop: 12
%% chuong: 4
%% bai: 13
%% dang: 12.13.15
%% loai: TN4
%% muc: TH
%% dap_an: "D"
%% nguon: "(NBV)-ĐỀ SỐ 16-MỖI NGÀY 1 ĐỀ THI 2025.pdf (D16, MỖI NGÀY 1 ĐỀ - Đề số 16), Phần I Câu 8"
%% kien_thuc: "Bài 12 (tích phân), Bài 13 (diện tích hình phẳng bằng tích phân)"
%% trang_thai: cho-duyet
%% ly_do: "Dạng mới đề xuất 12.13.15 chờ giáo viên duyệt"
%% ghi_chu: "Đáp án gốc D đúng."
\begin{ex}
Cho hàm số $y=f(x)$ liên tục trên $\mathbb R$ và đồ thị hàm số $y=f(x)$ cắt trục hoành tại đúng hai điểm phân biệt $x_1,x_2$ $(x_1<x_2)$. Gọi $S$ là diện tích hình phẳng giới hạn bởi đồ thị hàm số $y=f(x)$ và trục $Ox$. Biết rằng $f(x)\le0$ với mọi $x\in[x_1;x_2]$, giá trị của $\displaystyle\int_{x_1}^{x_2}f(x)\,\mathrm dx$ bằng
\choice{$(x_1-x_2)S$}{$S$}{$(x_2-x_1)S$}{\True $-S$}
\loigiai{$S=\displaystyle\int_{x_1}^{x_2}|f(x)|\,\mathrm dx=\int_{x_1}^{x_2}\bigl(-f(x)\bigr)\mathrm dx$ vì $f(x)\le0$ trên $[x_1;x_2]$, suy ra $\displaystyle\int_{x_1}^{x_2}f(x)\,\mathrm dx=-S$. Chọn D.}
\end{ex}
